\section{Finite certificate for the four-input precode}
\label{app:small-certificate}

We evaluate~\eqref{eq:small-union} using rational inequalities. All
notation refers to Appendix~\ref{sec:small-scans}, and binomial coefficients
outside their usual range are zero. Throughout, $\rho=e/(n-e)$.

\paragraph{Last-round tail.}
Decompose an output support into runs of nonzero outputs. A run is
\emph{closed} if a return ends it and \emph{open} if it reaches the last
position. For $h,j>0$, the number of support pairs with input weight $h$,
output weight $j$, and exactly $c$ returns is
\[
 \binom{n-j}{c}\left[
 \binom{j-1}{c-1}\binom{j-c}{h-2c}
 +\binom{j-1}{c}\binom{j-c-1}{h-2c-1}\right],
\]
and $A_g(h,j)$ weights this number by $g^c$. The two terms correspond to
a closed and an open final run; in the second term $c=0$ is allowed.
Summing over $j\le e$ and extending the number of available zero
positions gives
\begin{equation}\label{eq:small-tail}
 \frac{H_g(h,e)}{\binom nh}
 \le 2\sum_{c=0}^{\lfloor h/2\rfloor}
       \binom hc\binom{h-c}{c}g^c \rho^{h-c}
 \le 2(\rho+2\sqrt{g\rho})^h,\qquad h\le2e.
\end{equation}
For the first inequality, the closed-run binomial comparison is
\[
 \frac{\binom nc\binom e{h-c}}{\binom nh}
 =\binom hc\frac{\binom e{h-c}}{\binom{n-c}{h-c}}
 \le\binom hc \rho^{h-c},
\]
and the open case has the same bound with $c$ closed runs. The remaining
coefficient in each case is at most $\binom{h-c}{c}$. For the second
inequality, $h!/(c!^2(h-2c)!)=\binom hc\binom{h-c}{c}$, so the sum is
the constant coefficient in $u$ of the Laurent polynomial
$\bigl(\rho+\sqrt{g\rho}\,(u+u^{-1})\bigr)^h$. All its coefficients are
nonnegative, so this constant coefficient is at most its value at $u=1$,
namely $(\rho+2\sqrt{g\rho})^h$.

Fix a rational $t<1$ with $t>\rho$ and $(t-\rho)^2\ge4g\rho$, that is,
$\rho+2\sqrt{g\rho}\le t$. Then~\eqref{eq:small-tail} gives
$H_g(h,e)/\binom nh\le2t^h$ for $h\le2e$. Since $g$ depends on $\tau$,
so does the admissible range of $t$.

\paragraph{Bulk source weights.}
Choose positive rational $x,v,z$, rational $\tau\ge1$, $t$ as above, and
$0<y\le t$. Set
\[
 \lambda=\max\{1+xyv,\;gx/v+(1+x)y\},\qquad C=\max\{1,1/v\}.
\]
The transfer matrix in~\eqref{eq:small-transfer} maps $[1,v]^T$ to a
vector bounded coordinatewise by $\lambda[1,v]^T$. Its generating
function is consequently at most $C\lambda^n$. Combining coefficient
bounds with~\eqref{eq:small-tail}, the contribution of a source weight
$s$ to~\eqref{eq:small-union} is at most
\begin{equation}\label{eq:small-endpoint}
 \frac{2(n+1)C}{\tau} Q(z,\tau)^L\lambda^n
 (t/y)^{2e}(xz)^{-s}(s/n)^s(1-s/n)^{n-s}.
\end{equation}
The factor $(n+1)(s/n)^s(1-s/n)^{n-s}$ bounds $1/\binom ns$, by the
estimate $\binom ns\ge2^{nH_2(s/n)}/(n+1)$. With witnesses fixed, the
logarithm of~\eqref{eq:small-endpoint} is convex in $s$. Its values at
an interval's two endpoints therefore bound its value at every interior
integer. Multiplying by the number of integers in the interval bounds
the interval's total.

\paragraph{Sparse source weights.}
For $1\le s\le128$, retain the return weight in the sparse
calculation. With $\tau\ge1$, $g=1/q+(1-1/q)/\tau$, and $t$ as above,
put $b_g=t/(1-t)$ and define
\begin{equation}\label{eq:small-sparse-moment}
\begin{aligned}
 M_s(t;g)={}&\sum_{j=1}^{\lfloor s/2\rfloor}
 \binom nj\binom{s-j-1}{j-1}g^j b_g^{s-j}\\
 &+\sum_{j=1}^{\lceil s/2\rceil}
 \binom n{j-1}\binom{s-j}{j-1}g^{j-1}b_g^{s-j+1}.
\end{aligned}
\end{equation}
An open final run ends with no return, which explains the factor
$g^{j-1}$ in the second sum. To obtain the bound, extend each nonzero
run to any positive length. For $j$ closed runs the generating function
is $g^j x^{2j}[t/(1-t(1+x))]^j$. An open final run has one fewer input
boundary and one fewer return. Thus $M_s(t;g)$ bounds
$\sum_h A_g(s,h)t^h$. The entire sparse contribution is at most
\begin{equation}\label{eq:small-sparse-weighted}
 P_{\rm sparse}:=\frac2\tau\sum_{s=1}^{128}
 \frac{[z^s]Q(z,\tau)^L}{\binom ns}M_s(t;g).
\end{equation}
At $\tau=1$, $g=1$ and this is the unweighted sparse bound, which
ignores coefficient randomness. The value of $\tau$ is a proof parameter
and does not change the encoder coefficients.

\paragraph{Certified values.}
For the row $(q,e)=(32,209715)$, the unweighted sparse bound and
114 bulk intervals give failure below $2^{-24.85}$. For the implemented
alphabets, the weighted sparse bound with $\tau=\tau_{\rm sparse}$ gives
the following values. Bulk intervals use their own witnesses
in~\eqref{eq:small-endpoint}.
\[
\begin{array}{r@{\quad}r@{\quad}r@{\quad}r@{\quad}r}
q&e&\tau_{\rm sparse}&\text{intervals}&\log_2\Pr[d\le e]\text{ bound}\\\hline
40&230896&160&115&-44.11\\
64&251658&256&116&-45.65\\
128&277872&512&119&-47.87\\
256&288358&1024&118&-50.24\\
512&298844&2048&126&-52.46\\
1024&301989&4096&119&-54.72\\
2048&304087&8192&119&-56.90
\end{array}
\]
Each bulk interval contributes less than $2^{-80}$, and the intervals
together cover every integer from 129 through $4e$. These certificates
prove Theorem~\ref{thm:small-scans}. The configuration with the largest
certified distance in this four-input family uses $e=304087$, so its
bulk cover ends at 1216348. Exact sparse evaluation and outward
logarithm bounds give
\[
 \log_2(P_{\rm sparse}+119\cdot2^{-80})
 \le -56.907879669916.
\]
At $q=2048$ and $e=230896$, a separate cover of 115 intervals gives
the upper logarithm $-58.451178036656$. Rational entropy bounds also
verify
\[
 H_2(230897/2097152)>1/2,
\]
so the $q=40$ distance bound exceeds the binary rate-$1/2$
Gilbert--Varshamov benchmark, as stated after Theorem~\ref{thm:small-scans}.
The certificates and their verifiers are listed in
Appendix~\ref{app:reproduction}.
